%\ifodd\value{page}\hbox{}\newpage\fi
\chapter*{Śrī Lalitā Sahasranāma — Śloka 3}
\addcontentsline{toc}{chapter}{Śloka 3 - Nomi 6,7}

\begin{sloka}
{\sanskritfont
\begin{tabular}{c}
मनोरूपेक्षुकोदण्डा पञ्चतन्मात्रसायका ।\\
निजारुणप्रभापूरमज्जद्ब्रह्माण्डमण्डला ॥ ३ ॥
\end{tabular}

\vspace{1.5em}

{\middlesize \iastfont
manorūpekṣukodaṇḍā pañcatanmātrasāyakā |\\
nijāruṇaprabhāpūramajjadbrahmāṇḍamaṇḍalā ||3||}

\vspace{1em}

{\middlesize
\anvaya{%
मनोरूपेक्षुकोदण्डा पञ्चतन्मात्रसायका
निजारुणप्रभापूरमज्जद्ब्रह्माण्डमण्डला ।
}
}

\middlesize \iastfont \itmean{%
«Il cui arco di canna da zucchero ha la forma della mente,
le cui frecce sono i cinque \textit{tanmātra};
colei che immerge l’intero universo nella pienezza del proprio splendore rosso.»%
}

}
\end{sloka}

\wordmeaning{%
\wordline{मनः-रूप-इक्षु-कोदण्डा}{manaḥ-rūpa-ikṣu-kodaṇḍā}
  { che ha come arco di canna da zucchero la mente;}%
\wordline{पञ्च-तन्मात्र-सायका}{pañca-tanmātra-sāyakā}
  { che ha come frecce i cinque \textit{tanmātra} — \textit{śabda-tanmātra}: suono, \textit{sparśa-tanmātra}: tatto, \textit{rūpa-tanmātra}: forma/colore (visibile), \textit{rasa-tanmātra}: sapore, \textit{gandha-tanmātra}: odore;}%
\wordline{निज-अरुण-प्रभा-पूर-मज्जत्-ब्रह्माण्ड-मण्डला}
  {nija-aruṇa-prabhā-pūra-\\
  majjad-brahmāṇḍa-maṇḍalā}
  { colei che immerge l’intero universo nella pienezza del proprio splendore rosso;}%
}

Lo \textit{Śloka} 3 esplicita il principio della manifestazione mostrando che l’azione creativa di \textit{Lalitā} non avviene mediante strumenti materiali, ma attraverso strutture sottili e cognitive. L’arco, identificato con la \textit{manas} e descritto come \textit{ikṣu-kodaṇḍa}, indica che la proiezione della mente opera secondo una modalità di dolce attrazione e non di costrizione, rendendo l’esperienza desiderabile prima ancora che determinata. Le frecce, i \textit{pañca-tanmātra}, rappresentano le determinazioni elementari della percezione — suono, contatto, forma, sapore e odore — mediante le quali l’esperienza prende forma prima di diventare mondo sensibile. Presentati come frecce, i \textit{tanmātra} indicano che la manifestazione è un atto mirato e intelligibile, non un processo casuale o materiale. Il verso conclusivo afferma che l’intero universo è immerso nella \textit{aruṇa-prabhā} di \textit{Lalitā}, indicando che la creazione non si separa mai dalla coscienza che la genera, ma permane interamente contenuta nella pienezza del suo stesso splendore.

Lo \textit{Śloka} 3 presenta due nāma principali, che esplicitano il meccanismo della manifestazione. \textbf{\textit{Manorūpekṣukodaṇḍā}} indica che l’arco dell’azione divina è la mente stessa, principio di proiezione e configurazione; \textbf{\textit{pañca-tanmātra-sāyakā}} afferma che le frecce sono le cinque determinazioni sottili della percezione, attraverso cui l’esperienza prende forma. L’espressione \textit{nijāruṇa-prabhā-pūra-majjad-brahmāṇḍa-maṇḍalā }non introduce un ulteriore nome distinto, ma descrive l’effetto cosmico unitario di questi \textit{nāma}, affermando che l’intero universo rimane immerso nella pienezza dello splendore della Dea.